\documentclass[11pt]{article}

\usepackage[a4paper,margin=1in]{geometry}
\usepackage{amsmath,amssymb,mathtools}
\usepackage{booktabs,longtable,array}
\usepackage[hidelinks]{hyperref}

\newcommand{\R}{\mathbb{R}}
\newcommand{\dd}{\mathrm{d}}
\newcommand{\clip}{\operatorname{clip}}
\newcommand{\proj}{\operatorname{Proj}}

\title{Dynamic EV Charging-Station Pricing:\\Problem Statement, Mathematical Model, and Numerical Solution}
\author{}
\date{}

\begin{document}
\maketitle

\section{Problem statement}

We consider a directed transportation network containing ordinary road links
and charging-station links. Travellers belong to classes such as electric
vehicles (EVs) and non-electric vehicles (NEVs). For every origin--destination
(OD) pair, travellers choose among a bounded set of feasible routes. EV routes
contain one reachable charging station, whereas NEV routes contain only road
links.

The problem contains two coupled decision levels:

\begin{enumerate}
    \item \textbf{Inner traffic and route-choice equilibrium.} For fixed
    charging prices, road occupancies, station occupancies, and route flows
    evolve according to a system of ordinary differential equations (ODEs).
    Travellers gradually move flow toward routes with lower perceived cost.
    The inner problem is solved until the complete ODE is close to a fixed
    point.

    \item \textbf{Outer strategic pricing problem.} Every station changes its
    price to increase its own equilibrium profit. Changing one price changes
    travellers' route choices, congestion, station occupancy, and throughput;
    therefore profit must be evaluated after solving the inner equilibrium
    again.
\end{enumerate}

For a price vector $\boldsymbol{\psi}$, let
$\mathbf{z}^{\star}(\boldsymbol{\psi})$ denote the resulting inner
equilibrium. Station $s$ seeks to increase
\begin{equation}
    \Pi_s\!\left(\boldsymbol{\psi},
    \mathbf{z}^{\star}(\boldsymbol{\psi})\right).
\end{equation}
Thus the overall problem is a nested equilibrium problem rather than a single
static shortest-path calculation.

\section{Notation}

\begin{longtable}{>{\raggedright\arraybackslash}p{0.19\textwidth}
                  >{\raggedright\arraybackslash}p{0.73\textwidth}}
\toprule
Symbol & Meaning \\
\midrule
\endfirsthead
\toprule
Symbol & Meaning \\
\midrule
\endhead
$G=(V,E)$ & Directed transportation graph, with node set $V$ and link set $E$. \\
$E_r$ & Set of ordinary road links. \\
$E_s$ & Set of charging-station links. \\
$\mathcal{C}$ & Set of vehicle classes, for example $\{\mathrm{EV},\mathrm{NEV}\}$. \\
$\mathcal{W}$ & Set of OD pairs. \\
$w=(o_w,d_w)$ & OD pair with origin $o_w$ and destination $d_w$. \\
$\lambda_w(t)$ & Total exogenous demand rate for OD pair $w$ at time $t$. \\
$\theta_{wc}$ & Share of OD demand belonging to class $c$, with $\sum_c\theta_{wc}=1$. \\
$\lambda_{wc}(t)$ & Class-specific demand, $\lambda_{wc}(t)=\theta_{wc}\lambda_w(t)$. \\
$\mathcal{P}_{wc}$ & Bounded set of feasible paths for OD/class group $(w,c)$. \\
$p$ & A path in $\mathcal{P}_{wc}$. \\
$\delta_{ep}$ & Link--path incidence: $1$ if link $e$ belongs to path $p$, otherwise $0$. \\
$x_{ec}(t)$ & Occupancy/state of road link $e\in E_r$ for class $c$. \\
$x_s(t)$ & Occupancy/state of charging station $s\in E_s$. \\
$y_p(t)$ & Replicator route-flow state associated with path $p$. \\
$S_{wc}(t)$ & Current route-flow total, $S_{wc}(t)=\sum_{p\in\mathcal{P}_{wc}}y_p(t)$. \\
$q_p(t)$ & Physical demand assigned to path $p$. \\
$u_{ec}(t)$ & Arrival rate to road link $e$ for class $c$. \\
$u_s(t)$ & Arrival rate to charging station $s$. \\
$f_{ec}(t)$ & Road outflow rate from link $e$ for class $c$. \\
$\rho_s(t)$ & Charging-station throughput. \\
$a_e$ & Linear road discharge coefficient. \\
$L_e$ & Road flow/capacity scale used by the latency function. \\
$\ell_e^0$ & Road latency coefficient. \\
$a_s$ & Unsaturated station discharge coefficient. \\
$\mu_s$ & Maximum station service rate. \\
$K_s$ & Station saturation occupancy, $K_s=\mu_s/a_s$. \\
$W_s(x_s)$ & Station waiting-time function. \\
$\phi_s^0$ & Fixed non-price component of station cost. \\
$\alpha$ & Weight assigned to station waiting time. \\
$\gamma$ & Weight converting monetary price into perceived generalized cost. \\
$\psi_s$ & Price charged by station $s$. \\
$c_s$ & Per-vehicle operating cost of station $s$. \\
$C_e$ & Generalized cost of link $e$. \\
$C_p$ & Generalized cost of path $p$. \\
$\overline C_{wc}$ & Current flow-weighted average path cost for group $(w,c)$. \\
$\eta$ & Route-choice/replicator adaptation rate. \\
$\mathbf{x}$ & Vector containing every road and station state. \\
$\mathbf{y}$ & Vector containing every multi-route replicator state. \\
$\mathbf{z}$ & Complete inner state, $\mathbf{z}=[\mathbf{x}^{\mathsf T},\mathbf{y}^{\mathsf T}]^{\mathsf T}$. \\
$F$ & Complete ODE right-hand side, $\dot{\mathbf z}=F(\mathbf z;\boldsymbol\psi)$. \\
$R$ & Inner equilibrium residual, $R=\|\dot{\mathbf{x}}\|_2+\|\dot{\mathbf{y}}\|_2$. \\
$\varepsilon$ & Inner stopping tolerance; the implementation uses $10^{-6}$. \\
$T_{\max}$ & Maximum inner integration horizon; the default is $1000$. \\
$\Pi_s$ & Profit rate of station $s$. \\
$\Delta$ & Finite-difference price perturbation; the default is $0.02$. \\
$g_s$ & Estimated derivative of station $s$'s profit with respect to its price. \\
$g_{\max}$ & Gradient clipping bound; the default is $5$. \\
$\kappa$ & Outer price adaptation gain; the default is $0.1$. \\
$\Delta\tau$ & Outer-loop step duration; the default is $1$. \\
$\psi_{\max}$ & Maximum permitted station price; the default is $3$. \\
$k$ & Outer pricing-iteration index. \\
\bottomrule
\end{longtable}

\section{Step 1: feasible-route construction}

For each group $(w,c)$, the implementation first removes links that cannot be
used by class $c$. It then constructs a bounded set $\mathcal{P}_{wc}$ of
short, simple, feasible alternatives. NEV paths contain no station link. EV
paths contain exactly one station link and are formed from a road path from the
origin to a station, the station link itself, and a road path from the station
to the destination.

Bounding $|\mathcal{P}_{wc}|$ is essential because enumerating every simple
path can grow exponentially with network size. In preview mode, one best road
prefix and suffix is retained per reachable station, followed by a limit of
eight paths per OD/class group. Class-filtered road graphs and repeated path
segments are cached during construction. The route set is part of the model:
changing it can change the computed equilibrium.

\section{Step 2: class-specific demand and route flow}

The demand of class $c$ for OD pair $w$ is
\begin{equation}
    \lambda_{wc}(t)=\theta_{wc}\lambda_w(t).
\end{equation}
For a group with multiple paths, the physical flow assigned to path $p$ is
\begin{equation}
    q_p(t)=\lambda_{wc}(t)\frac{y_p(t)}{S_{wc}(t)},
    \qquad
    S_{wc}(t)=\sum_{r\in\mathcal{P}_{wc}}y_r(t).
    \label{eq:path-flow}
\end{equation}
If there is only one feasible path, all class demand is assigned to it. If
$S_{wc}$ is numerically zero, equal route fractions are used defensively.
Because the solver projects route flows onto
\begin{equation}
    \mathcal{Y}_{wc}=
    \left\{\mathbf y_{wc}\geq 0:
    \sum_{p\in\mathcal{P}_{wc}}y_p=\lambda_{wc}\right\},
\end{equation}
normally $S_{wc}=\lambda_{wc}$ and hence $q_p=y_p$.

\section{Step 3: road and station dynamics}

For road link $e$ and class $c$, the outflow is
\begin{equation}
    f_{ec}(t)=a_e x_{ec}(t).
\end{equation}
The road conservation equation is
\begin{equation}
    \dot{x}_{ec}(t)=u_{ec}(t)-f_{ec}(t).
    \label{eq:road-dynamics}
\end{equation}

For charging station $s$, throughput is saturated:
\begin{equation}
    \rho_s(x_s)=\min\{a_sx_s,\mu_s\},
    \label{eq:throughput}
\end{equation}
and station occupancy follows
\begin{equation}
    \dot{x}_s(t)=u_s(t)-\rho_s(x_s(t)).
    \label{eq:station-dynamics}
\end{equation}

At an origin, the incoming flow to the first link of every route is $q_p$.
At an internal link, the outgoing flow is divided among successor links using
the route-flow proportions. If $e$ is followed by $j$ for some class-$c$
paths, the implemented turning fraction can be written as
\begin{equation}
    \tau_{e\rightarrow j,c}
    =\frac{\sum_{p:\,(e,j)\in p,\,p\in\mathcal{P}_{wc}}q_p}
    {\sum_{p:\,e\in p,\,p\in\mathcal{P}_{wc}}q_p}.
\end{equation}
Consequently, $f_{ec}\tau_{e\rightarrow j,c}$ contributes to the arrival
rate of successor $j$. This construction conserves flow through the network.

\section{Step 4: generalized link and path costs}

Let total occupancy on road $e$ be
\begin{equation}
    x_e^{\mathrm{tot}}=\sum_{c\in\mathcal C}x_{ec},
\end{equation}
and let total road outflow be $f_e=a_ex_e^{\mathrm{tot}}$. The implemented
road cost is
\begin{equation}
    C_e^{r}(x_e^{\mathrm{tot}})
    =\ell_e^0
    \frac{f_e/L_e}{1-f_e/L_e},
    \qquad 0\leq \frac{f_e}{L_e}<1.
    \label{eq:road-cost}
\end{equation}
Numerically, $f_e/L_e$ is clipped below $1$ to avoid division by zero.

The saturation occupancy and waiting time at station $s$ are
\begin{align}
    K_s &= \frac{\mu_s}{a_s},\\
    W_s(x_s)&=\frac{\max\{x_s-K_s,0\}}{\mu_s}.
\end{align}
The station's generalized user cost is
\begin{equation}
    C_s^{\mathrm{ch}}(x_s,\psi_s)
    =\phi_s^0+\alpha W_s(x_s)+\gamma\psi_s.
    \label{eq:station-cost}
\end{equation}
The total cost of path $p$ is additive:
\begin{equation}
    C_p(\mathbf x,\boldsymbol\psi)
    =\sum_{e\in p\cap E_r}C_e^r
    +\sum_{s\in p\cap E_s}C_s^{\mathrm{ch}}.
    \label{eq:path-cost}
\end{equation}

\section{Step 5: route-choice dynamics}

For each multi-route group $(w,c)$, define its current flow-weighted average
cost as
\begin{equation}
    \overline C_{wc}
    =\frac{\sum_{p\in\mathcal P_{wc}}y_pC_p}
    {S_{wc}},
    \qquad S_{wc}=\sum_{p\in\mathcal P_{wc}}y_p.
    \label{eq:average-cost}
\end{equation}
The route-flow state evolves according to replicator dynamics:
\begin{equation}
    \dot y_p
    =\eta y_p\left(\overline C_{wc}-C_p\right),
    \qquad p\in\mathcal P_{wc}.
    \label{eq:replicator}
\end{equation}
If $C_p<\overline C_{wc}$, then $\dot y_p>0$ and flow moves toward path $p$.
If $C_p>\overline C_{wc}$, then $\dot y_p<0$.

The current-total denominator in Equation~\eqref{eq:average-cost} preserves
the group total:
\begin{align}
    \frac{\dd S_{wc}}{\dd t}
    &=\sum_p\dot y_p\\
    &=\eta\left(
       \overline C_{wc}\sum_p y_p-\sum_p y_pC_p
      \right)\\
    &=\eta\left(
       \frac{\sum_p y_pC_p}{S_{wc}}S_{wc}-\sum_p y_pC_p
      \right)=0.
\end{align}
This identity is the reason for normalizing by the current sum $S_{wc}$ rather
than by the target demand. It prevents a small numerical deviation in the
route-flow total from being amplified by the vector field.

At a route-choice equilibrium, every used path has the same minimum cost:
\begin{equation}
    y_p^{\star}>0
    \quad\Longrightarrow\quad
    C_p(\mathbf x^{\star},\boldsymbol\psi)
    =\overline C_{wc}^{\star}.
    \label{eq:wardrop}
\end{equation}
This is the Wardrop-type user-equilibrium interpretation of the replicator
fixed point.

\section{Step 6: inner equilibrium computation}

Collect all states into
\begin{equation}
    \mathbf z=
    \begin{bmatrix}\mathbf x\\\mathbf y\end{bmatrix},
    \qquad
    \dot{\mathbf z}=F(\mathbf z;\boldsymbol\psi).
\end{equation}
For fixed prices, an exact equilibrium satisfies
\begin{equation}
    F(\mathbf z^{\star};\boldsymbol\psi)=\mathbf 0.
\end{equation}
The implementation accepts a numerical equilibrium when
\begin{equation}
    R(\mathbf z;\boldsymbol\psi)
    =\left\|\dot{\mathbf x}\right\|_2
     +\left\|\dot{\mathbf y}\right\|_2
    \leq\varepsilon,
    \qquad \varepsilon=10^{-6}.
    \label{eq:residual}
\end{equation}

The ODE is integrated with SciPy's implicit BDF method because the coupled
traffic and replicator dynamics can be stiff. The default numerical settings
are relative tolerance $10^{-7}$, absolute tolerance $10^{-9}$, maximum BDF
step $15$, chunk length $50$, and maximum horizon $T_{\max}=1000$.

For equilibrium integration only, the route-choice derivative is multiplied
by a positive time-scale factor $20$. This accelerates slow adaptation caused
by the small price perturbations used for finite differences. Because
$20\dot{\mathbf y}=\mathbf 0$ if and only if
$\dot{\mathbf y}=\mathbf 0$, the equilibrium points are unchanged. The
stopping residual in Equation~\eqref{eq:residual} is evaluated using the
original unscaled dynamics. Ordinary time-trajectory simulation also retains
the original configured value of $\eta$.

After each chunk, route flows are projected back onto their simplex. Define
$\widetilde y_p=\max\{y_p,0\}$. If
$\sum_r\widetilde y_r>0$, the projection is
\begin{equation}
    y_p\leftarrow
    \lambda_{wc}\frac{\widetilde y_p}
    {\sum_{r\in\mathcal P_{wc}}\widetilde y_r}.
    \label{eq:projection}
\end{equation}
If every clipped value is zero, demand is distributed equally. The residual
in Equation~\eqref{eq:residual} is recomputed after projection; therefore an
event detected on an unprojected trajectory is not accepted blindly.

The solver also uses a sparse Jacobian pattern, warm-starts nearby price
profiles from their previous equilibrium, caches verified equilibria, and
returns a non-convergence warning rather than representing a finite endpoint
as an exact fixed point.

\section{Step 7: station profit and strategic price update}

At inner equilibrium, station $s$ earns
\begin{equation}
    \Pi_s(\boldsymbol\psi)
    =(\psi_s-c_s)\rho_s
    \left(x_s^{\star}(\boldsymbol\psi)\right).
    \label{eq:profit}
\end{equation}
The dependence of equilibrium throughput on every station price makes an
analytic derivative inconvenient. The implementation uses a coordinate-wise
central finite difference. For station $s$ at iteration $k$, define
\begin{align}
    \psi_{s,k}^{+}&=\min\{\psi_{s,k}+\Delta,\psi_{\max}\},\\
    \psi_{s,k}^{-}&=\max\{\psi_{s,k}-\Delta,0\}.
\end{align}
All other station prices remain fixed. After separately solving the inner
equilibrium at the two perturbed price vectors,
\begin{equation}
    g_{s,k}
    \approx
    \frac{\Pi_s(\boldsymbol\psi_k^{+})
          -\Pi_s(\boldsymbol\psi_k^{-})}
         {\psi_{s,k}^{+}-\psi_{s,k}^{-}}.
    \label{eq:finite-difference}
\end{equation}
The estimate is clipped to $[-g_{\max},g_{\max}]$, and projected gradient
ascent gives
\begin{equation}
    \psi_{s,k+1}
    =\proj_{[0,\psi_{\max}]}
      \left(\psi_{s,k}+\Delta\tau\,\kappa\,
      \clip(g_{s,k},-g_{\max},g_{\max})\right).
    \label{eq:price-update}
\end{equation}
The implementation allows at most $N=30$ outer iterations by default and can
stop earlier after the requested number of consecutive stable, fully
converged iterations. With $m$ stations, the upper bound on inner solves is
\begin{equation}
    N(1+2m)+1
\end{equation}
inner solves: one nominal and $2m$ perturbed solves per iteration, followed by
one final solve. Independent perturbation solves are evaluated in parallel
when worker processes are available.

\section{Worked numerical example}

Consider one EV OD pair with constant demand $\lambda=1$ vehicle/time and two
paths, $p_1$ through station $s_1$ and $p_2$ through station $s_2$. Assume
\begin{equation}
    y_1=0.6,\qquad y_2=0.4,\qquad S=y_1+y_2=1.
\end{equation}
From Equation~\eqref{eq:path-flow},
\begin{equation}
    q_1=1\frac{0.6}{1}=0.6,
    \qquad
    q_2=1\frac{0.4}{1}=0.4.
\end{equation}

\subsection*{Road costs}

Let both roads have $\ell^0=0.2$, $L=4$, and $a=1$. Let their occupancies be
$x_1=1$ and $x_2=0.5$. For road 1,
\begin{align}
    f_1&=a x_1=1,\\
    C_1^r&=0.2\frac{1/4}{1-1/4}
          =0.066667.
\end{align}
For road 2,
\begin{align}
    f_2&=a x_2=0.5,\\
    C_2^r&=0.2\frac{0.5/4}{1-0.5/4}
          =0.028571.
\end{align}

\subsection*{Station throughput, waiting time, and cost}

Let both stations have $a_s=0.5$, $\mu_s=2$, $\phi_s^0=0.1$,
$\alpha=0.3$, and $\gamma=1$. Thus
\begin{equation}
    K_s=\frac{2}{0.5}=4.
\end{equation}
Suppose $x_{s_1}=5$, $\psi_{s_1}=0.5$, $x_{s_2}=2$, and
$\psi_{s_2}=0.7$. Then
\begin{align}
    \rho_{s_1}&=\min\{0.5(5),2\}=2,\\
    W_{s_1}&=\frac{\max\{5-4,0\}}{2}=0.5,\\
    C_{s_1}^{\mathrm{ch}}&=0.1+0.3(0.5)+1(0.5)=0.75,
\end{align}
whereas
\begin{align}
    \rho_{s_2}&=\min\{0.5(2),2\}=1,\\
    W_{s_2}&=0,\\
    C_{s_2}^{\mathrm{ch}}&=0.1+0+1(0.7)=0.8.
\end{align}

\subsection*{Path costs and route adaptation}

Each path contains its corresponding road and station, so
\begin{align}
    C_{p_1}&=0.066667+0.75=0.816667,\\
    C_{p_2}&=0.028571+0.8=0.828571.
\end{align}
The current average cost is
\begin{align}
    \overline C
    &=0.6(0.816667)+0.4(0.828571)\\
    &=0.821429.
\end{align}
With $\eta=0.05$, Equation~\eqref{eq:replicator} gives
\begin{align}
    \dot y_1
    &=0.05(0.6)(0.821429-0.816667)
      =1.4286\times10^{-4},\\
    \dot y_2
    &=0.05(0.4)(0.821429-0.828571)
      =-1.4284\times10^{-4}.
\end{align}
The cheaper path gains flow and the more expensive path loses flow. Apart from
rounding,
\begin{equation}
    \dot y_1+\dot y_2=0,
\end{equation}
so total OD/class demand is preserved.

\subsection*{Physical-state derivative and equilibrium residual}

Assume each path has a unique road followed by its station. Road arrivals are
$0.6$ and $0.4$, while road outflows are $1$ and $0.5$. Therefore
\begin{align}
    \dot x_1&=0.6-1=-0.4,\\
    \dot x_2&=0.4-0.5=-0.1.
\end{align}
The road outflows become station arrivals, so
\begin{align}
    \dot x_{s_1}&=1-2=-1,\\
    \dot x_{s_2}&=0.5-1=-0.5.
\end{align}
Thus
\begin{align}
    \|\dot{\mathbf x}\|_2
    &=\sqrt{(-0.4)^2+(-0.1)^2+(-1)^2+(-0.5)^2}\\
    &=\sqrt{1.42}=1.19164,\\
    \|\dot{\mathbf y}\|_2
    &\approx\sqrt{(1.4286\times10^{-4})^2
                   +(-1.4284\times10^{-4})^2}\\
    &\approx2.020\times10^{-4},\\
    R&\approx1.19184.
\end{align}
Since $R\gg10^{-6}$, this state is not an equilibrium and BDF continues
integrating. At the accepted equilibrium, road and station inflows balance
their outflows and used path costs are equal, making both derivative norms
smaller than the required combined tolerance.

\subsection*{Profit and one price update}

Let operating cost be $c_{s_1}=0.2$. At the displayed state,
\begin{equation}
    \Pi_{s_1}=(0.5-0.2)(2)=0.6.
\end{equation}
For the actual gradient, the inner equilibrium must be solved again at both
perturbed prices. As an illustrative finite-difference step, let $\Delta=0.02$
and suppose those two equilibrium solves produce
\begin{align}
    \rho_{s_1}^{+}&=1.9 &&\text{at }\psi_{s_1}^{+}=0.52,\\
    \rho_{s_1}^{-}&=2.1 &&\text{at }\psi_{s_1}^{-}=0.48.
\end{align}
Then
\begin{align}
    \Pi_{s_1}^{+}&=(0.52-0.2)(1.9)=0.608,\\
    \Pi_{s_1}^{-}&=(0.48-0.2)(2.1)=0.588,\\
    g_{s_1}&\approx\frac{0.608-0.588}{0.52-0.48}=0.5.
\end{align}
Using $\kappa=0.1$ and $\Delta\tau=1$,
\begin{equation}
    \psi_{s_1,\mathrm{new}}
    =\proj_{[0,3]}\bigl(0.5+1(0.1)(0.5)\bigr)
    =0.55.
\end{equation}
Thus the station raises its price because its estimated equilibrium-profit
gradient is positive.

\section{Algorithm summary}

\begin{enumerate}
    \item Read the directed roads, stations, OD demands, vehicle classes, and
    model parameters.
    \item Generate a bounded feasible-route set for every OD/class group.
    \item Initialize road and station states to zero and distribute each
    multi-route demand equally across its paths.
    \item Fix the current station-price vector $\boldsymbol\psi$.
    \item Integrate the coupled traffic, station, and replicator ODE with BDF.
    After every chunk, project route flows onto their demand simplex.
    \item Stop when
    $\|\dot{\mathbf x}\|_2+\|\dot{\mathbf y}\|_2\leq10^{-6}$ or report that
    $T_{\max}$ was reached without full convergence.
    \item Compute station throughput and profit at the returned equilibrium.
    \item For every station, perturb its price upward and downward and repeat
    the inner solve to estimate the coordinate profit gradient.
    \item Apply the clipped, projected gradient-ascent price update.
    \item Repeat until the requested outer-step limit is reached or prices
    remain stable for the requested number of converged iterations, then
    perform one final equilibrium solve at the final prices.
\end{enumerate}

\section{What the final result means}

The output is an approximate coupled equilibrium. On the user side, no used
route has an incentive to move flow to a more expensive route. On the station
side, the final price vector is the endpoint of projected profit-gradient
updates and is interpreted as an approximate local strategic-price
equilibrium when price changes are small. A result is numerically trustworthy
only when all required inner solves satisfy the residual tolerance, route-flow
conservation error is negligible, and no important route set was truncated.

\end{document}
